\documentclass[10pt,a4paper]{article}
\usepackage[utf8]{inputenc}
\usepackage{amsmath}
\usepackage{amsfonts}
\usepackage{amssymb}
\usepackage{geometry}
\usepackage{multicol}
\usepackage{enumitem}
\usepackage{xcolor}

\geometry{margin=0.6in, top=0.5in, bottom=0.5in}

\newcommand{\redflag}[1]{\textcolor{red}{\textbf{[!] #1}}}
\newcommand{\greennote}[1]{\textcolor{green!50!black}{\textbf{[CORRECT] #1}}}
\newcommand{\examfocus}[1]{\textcolor{orange!80!black}{\textbf{[FOCUS] #1}}}
\newcommand{\trap}[1]{\textcolor{purple}{\textbf{[TRAP] #1}}}
\newcommand{\keyterm}[1]{\textcolor{red!80!black}{\textbf{[KEY] #1}}}

\setlength{\parindent}{0pt}
\setlength{\parskip}{0.1em}
\setlist{nosep, left=0pt, itemsep=0.05em, parsep=0em, topsep=0em}

\begin{document}

\centering
{\Large \textbf{CAMS Tutorial: High Net Worth Individuals}} \\
{\large Hard Topics – Private Banking, Trusts, PEPs, Offshore Risks} \\
\vspace{0.2cm}
\rule{\textwidth}{0.5pt}
\vspace{0.2cm}

\tiny
\raggedright

\section*{How Hard HNW Questions Are Framed}
In the CAMS exam, HNW private banking questions test your ability to identify:
\begin{itemize}
\item \textbf{Inherent conflicts of interest} (RMs paid on AUM)
\item \textbf{Opaque structures} (trusts, shell companies, SPVs, PICs)
\item \textbf{PEP risks} (foreign, domestic, international org, family, close associates)
\item \textbf{Source of Wealth vs Source of Funds} (both must be verified)
\item \textbf{Offshore financial center red flags} (round tripping, rapid transfers, no economic purpose)
\item \textbf{High-value asset laundering} (art, yachts, crypto, insurance policies)
\item \textbf{Correspondent banking risks} (nested accounts, refusal to identify UBOs)
\item \textbf{Culture of compliance failures} (senior management overrides)
\end{itemize}

\section*{Key Ideas \& Explanations (No Questions)}

\subsection*{1. Conflict of Interest in Private Banking}
\keyterm{Inherent risk:} Relationship managers compensated on Assets Under Management (AUM) or revenue may ignore AML red flags to keep high-fee clients. \\
\examfocus{Performance compensation structures create ML risk.} \\
\trap{“The problem is lack of technology or process” – No, it’s human conflict of interest.} \\
\redflag{RM pushing to onboard PEP quickly without EDD = red flag.}

\subsection*{2. Complex Ownership Structures (Trusts + Offshore Companies)}
\keyterm{UBO obscurity:} Multiple layers (trust → BVI → Cyprus → Switzerland) hide ultimate beneficial owner. \\
\examfocus{Each offshore layer reduces transparency.} \\
\trap{“Trusts are legitimate wealth management tools” – They are, but can also be abused.} \\
\redflag{Client refuses to disclose beneficiaries or source of funds + multiple secrecy jurisdictions = maximum risk.}

\subsection*{3. Source of Wealth vs Source of Funds}
\keyterm{SoW:} How client accumulated total wealth over time (e.g., inheritance, business). \\
\keyterm{SoF:} Specific origin of money for this transaction (e.g., crypto sale). \\
\examfocus{BOTH must be independently verified.} \\
\trap{“Verifying SoW is sufficient for HNW clients” – Wrong. SoF also required.} \\
\redflag{Unverifiable crypto sale through privacy coin (Monero) + unregulated exchange = EDD and potential SAR.}

\subsection*{4. Trust Structure Red Flags}
\keyterm{Trustee self-dealing:} Son as trustee making loans to himself, buying yacht in trust name. \\
\keyterm{PEP settlor:} Once a PEP, always a PEP – even after death for risk assessment. \\
\examfocus{EDD required for foreign PEP settlors regardless of death.} \\
\trap{“Settlor is deceased, PEP status no longer matters” – Incorrect.} \\
\redflag{Unexplained wires to high-risk jurisdiction + crypto via unregulated VASP + refusal to explain = SAR.}

\subsection*{5. Offshore Financial Center Red Flags}
\keyterm{Round tripping:} Funds move Cayman → BVI → Bermuda → back to Cayman with no economic purpose. \\
\keyterm{Rapid transfers between offshore accounts.} \\
\keyterm{Shell company added as signatory.} \\
\examfocus{Each red flag requires EDD.} \\
\trap{“Tax planning using OFCs is legitimate; no red flags” – Wrong. Pattern matters.} \\
\redflag{Round tripping + rapid transfers without economic purpose = suspicious.}

\subsection*{6. Private Investment Companies (PICs)}
\keyterm{PICs in offshore jurisdictions reduce transparency.} \\
\examfocus{Client refusal to disclose underlying assets or investment strategy = red flag.} \\
\trap{“PICs are legitimate investment vehicles; no additional risk” – Not true when combined with opacity.} \\
\redflag{PICs + refusal to disclose + incoming wires from unrelated offshore entities + crypto = high risk.}

\subsection*{7. Trust Due Diligence (FATF)}
\keyterm{Must identify:} settlor, all trustees, protector, AND all beneficiaries (even discretionary). \\
\examfocus{Refusal to provide beneficiary names = grounds to reject account.} \\
\trap{“Trusts are private; beneficiaries do not need to be identified” – False under FATF.} \\
\redflag{All trust beneficiaries must be identified, including discretionary class.}

\subsection*{8. Special Purpose Vehicles (SPVs)}
\keyterm{SPVs can disguise true UBO.} \\
\examfocus{Rapid creation and dissolution of SPVs (each 12-18 months) + no clear business purpose = red flag.} \\
\trap{“SPVs are standard for holding single assets; no ML risk” – Pattern matters.} \\
\redflag{Rapid creation/dissolution + crypto + no business purpose = EDD required.}

\subsection*{9. Sovereign Wealth Funds (SWFs)}
\keyterm{SWFs are NOT automatically exempt from AML.} \\
\examfocus{EDD required for PEP board members; source of funds must be verified; Grey List status matters.} \\
\trap{“SWFs are government entities exempt from AML” – Incorrect.} \\
\redflag{SWF refusing to provide source of funds + PEP board members + Grey List = high risk.}

\subsection*{10. High Value Asset Red Flags}
\keyterm{Indicators of ML:} funds from secrecy jurisdiction, cash structuring, shell company ownership, unregulated dealer, income inconsistency, rapid resale at loss. \\
\examfocus{Criminals accept loss to clean money.} \\
\trap{“HNW individuals buy luxury assets regularly; no red flags” – Look for pattern.} \\
\redflag{Rapid resale at loss = money laundering (loss is laundering cost).}

\subsection*{11. Family Office Risks}
\keyterm{PEP patriarch (once a PEP, always a PEP).} \\
\keyterm{Multi-jurisdictional structure (Delaware registration, Swiss management) reduces transparency.} \\
\keyterm{Unidentified third parties (external advisor).} \\
\examfocus{Unexplained incoming wires from BVI shell company = investigate.} \\
\trap{“Patriarch left office 8 years ago, no longer PEP” – Once a PEP, always a PEP for risk.} \\
\redflag{Unidentified third parties in family office = red flag.}

\subsection*{12. Insurance Products ML Technique}
\keyterm{Policy loan layering:} (1) Purchase policy with illicit funds, (2) Take policy loan (clean proceeds), (3) Surrender policy for additional clean funds. \\
\examfocus{Repeating pattern with multiple insurers = systematic laundering.} \\
\trap{“Policy loans are legitimate; no ML risk” – Pattern reveals layering.} \\
\redflag{Policy loans from policies bought with illicit funds = layering.}

\subsection*{13. Cryptocurrency \& HNW Clients}
\keyterm{Red flags:} declining business revenues but large crypto purchases; VASP with weak KYC; DeFi protocol opacity; refusal to provide wallet addresses; unexplained rapid profits. \\
\examfocus{Declining income + large crypto = potential source of funds issue.} \\
\trap{“Crypto investing is normal for HNW” – Normal, but pattern matters.} \\
\redflag{Refusal to provide transaction history from DeFi = EDD required.}

\subsection*{14. Art and Collectibles Laundering}
\keyterm{Indicators:} freeport storage (no physical possession), rapid resale at loss, single-purpose LLCs for each transaction, different bank accounts for purchase and sale, refusal to explain strategy. \\
\examfocus{Rapid resale at loss + freeport = money laundering.} \\
\trap{“Art is a legitimate investment for HNW” – Yes, but pattern reveals laundering.} \\
\redflag{Loss accepted to clean money = laundering cost.}

\subsection*{15. Secured Loan Risks}
\keyterm{Money laundering via secured loans:} illicit funds used as collateral (portfolio bought with dirty money), loan proceeds are clean, early repayment after selling collateral creates legitimate paper trail. \\
\examfocus{Wire to high-risk jurisdiction + business with no record = red flag.} \\
\trap{“Secured loans are low risk; collateral guarantees repayment” – Collateral itself may be dirty.} \\
\redflag{Secured loans can launder funds if collateral purchased with illicit money.}

\subsection*{16. Multi-Jurisdiction Tax Evasion}
\keyterm{Tax evasion is a predicate offense for money laundering in most jurisdictions.} \\
\examfocus{CRS reporting mandatory regardless of client consent. Refusal to sign CRS form = red flag.} \\
\trap{“Tax evasion is not a bank’s concern” – It is a predicate offense.} \\
\redflag{Unexplained cross-border movements + income inconsistency = EDD and potential SAR.}

\subsection*{17. Shell Company Networks}
\keyterm{Strongest red flag:} massive discrepancy between declared revenue (\$5M) and transaction flow (\$500M). \\
\examfocus{Circular transactions = layering. Trust overlay adds final secrecy layer.} \\
\trap{“Shell companies are legal business structures; no risk” – Discrepancy reveals laundering.} \\
\redflag{Revenue vs transaction flow discrepancy = strongest red flag.}

\subsection*{18. Round Tripping}
\keyterm{Round tripping with consistent loss of 2-5\% = layering. Criminals accept loss to clean money.} \\
\examfocus{No legitimate business purpose + inability to explain = SAR required.} \\
\trap{“This could be legitimate treasury management” – Loss pattern indicates laundering.} \\
\redflag{Round tripping with consistent losses = money laundering (loss is fee).}

\subsection*{19. Foreign PEP Enhanced Due Diligence}
\keyterm{FATF requires EDD for foreign PEPs regardless of time since leaving office.} \\
\examfocus{Mathematical inconsistency (e.g., \$25M deposit vs \$12.5M max possible income) requires SAR if not resolved.} \\
\trap{“PEP status expires after leaving office” – False for foreign PEPs under FATF.} \\
\redflag{Foreign PEPs always require EDD, even years after leaving office.}

\subsection*{20. Domestic PEP Risk Assessment}
\keyterm{FATF recommends risk-based approach for domestic PEPs, not automatic EDD.} \\
\examfocus{Red flags (unexplained wealth, son’s government contracts with father’s involvement) require EDD regardless.} \\
\trap{“Domestic PEPs are exempt from EDD under FATF” – Not if red flags exist.} \\
\redflag{Unexplained wealth gap + conflict of interest = EDD required.}

\subsection*{21. International Organization PEPs (e.g., World Bank)}
\keyterm{FATF recommends risk-based assessment.} \\
\examfocus{Deposit pattern inconsistency (deposits stopped during World Bank employment and resumed after) is a red flag.} \\
\trap{“International organization PEPs are exempt from EDD” – Not under risk-based approach.} \\
\redflag{Inconsistent wealth accumulation patterns require EDD.}

\subsection*{22. Immediate Family of PEPs}
\keyterm{FATF definition includes immediate family members (spouse, children). Spouse is PEP by association.} \\
\examfocus{\$10M gift from \$150k salary is mathematically impossible → SAR required.} \\
\trap{“PEP definition applies only to the individual, not family” – False under FATF.} \\
\redflag{Immediate family of PEPs are also PEPs. Gifts from PEPs require source of funds verification.}

\subsection*{23. Close Associates of PEPs}
\keyterm{FATF definition includes close associates (business partners, friends).} \\
\examfocus{Hidden ownership through trust, allegations of contract steering, refusal to identify owners = EDD and potential SAR.} \\
\trap{“Close associates are not PEPs; only the individual is” – False.} \\
\redflag{Close associates of PEPs are considered PEPs for AML purposes.}

\subsection*{24. Correspondent Banking (FATF Rec. 13)}
\keyterm{EDD required for correspondent relationships.} \\
\examfocus{Red flags: respondent refuses to identify underlying clients (including PEPs), AML deficiencies, no independent audit, offshore jurisdiction. Termination or restriction required.} \\
\trap{“Correspondent banking risks can be managed with monitoring” – Not when respondent refuses UBO info.} \\
\redflag{Correspondent bank that refuses to identify underlying private banking clients should be terminated.}

\subsection*{25. Nested Accounts in Correspondent Banking}
\keyterm{Unauthorized nesting is a major risk. Correspondent bank has no visibility into underlying clients.} \\
\examfocus{PEP from sanctioned country using nested account = SAR required. Respondent’s refusal to terminate nesting = terminate correspondent relationship.} \\
\trap{“Nesting is permitted if respondent bank has adequate controls” – Unauthorized nesting is never acceptable.} \\
\redflag{Unauthorized nesting of private banking clients through correspondent accounts = terminate relationship.}

\subsection*{26. Exiting High-Risk Clients}
\keyterm{Banks can and should exit high-risk clients even with outstanding loans or letters of credit.} \\
\examfocus{Exit must be managed to avoid breach of contract. Revenue cannot be reason to retain high-risk client. SAR must be filed.} \\
\trap{“Banks cannot exit clients with outstanding credit facilities” – They can, carefully.} \\
\redflag{Revenue and outstanding loans are not valid reasons to retain a high-risk client.}

\subsection*{27. Culture of Compliance}
\keyterm{Weak culture indicators:} senior management overrides compliance recommendations for revenue; Board never reviews AML metrics; compensation 100\% revenue-based with no risk adjustment. \\
\examfocus{Senior management overriding compliance = critical failure.} \\
\trap{“Senior management can override compliance for commercial reasons” – No, that is a failure.} \\
\redflag{Senior management overriding compliance = culture of compliance failure.}

\subsection*{28. Estate Planning Trust Abuses}
\keyterm{Red flags:} settlor rapidly resigns as trustee, distributions to son personally (not in trustee capacity), son has no independent source of wealth, son moves to non-extradition country. \\
\examfocus{Trust distributions to individual not in trustee capacity = potential fraud or asset movement.} \\
\trap{“Trust distributions to beneficiaries are normal” – Not when done in personal capacity.} \\
\redflag{SAR required for suspected asset dissipation or fraud.}

\subsection*{29. Succession Planning Red Flags}
\keyterm{Migration to increasingly secretive jurisdictions (Cook Islands → Marshall Islands).} \\
\keyterm{Self-dealing (beneficiary becomes trustee of new trust).} \\
\keyterm{Relocation to non-disclosure country.} \\
\keyterm{Foundation funds diverted to secrecy jurisdiction with no charitable activity.} \\
\examfocus{Charitable funds diverted to secrecy jurisdiction = loss of charitable purpose + potential ML.} \\
\trap{“Succession planning is private; no bank oversight” – Red flags trigger bank obligations.} \\
\redflag{Foundation funds moved to Seychelles with no charitable activity = SAR.}

\subsection*{30. Offshore Trust Highest Risk Features}
\keyterm{Anonymous beneficiaries (no CDD possible).} \\
\keyterm{Flee clause – allows rapid asset movement to avoid freezing.} \\
\keyterm{Letter of wishes – secret instructions not subject to review.} \\
\keyterm{Pattern of trust creation/dissolution (each 18-24 months – moving assets to avoid detection).} \\
\keyterm{Art in freeport (opaque valuation and ownership).} \\
\keyterm{Crypto with unregulated VASP.} \\
\examfocus{High-risk structure requiring EDD and likely SAR.} \\
\trap{“Offshore trusts are standard for HNW international families” – Features above make them high risk.} \\
\redflag{Flee clauses + anonymous beneficiaries + rapid trust dissolution = highest risk.}

\vspace{0.5cm}
\rule{\textwidth}{0.5pt}
\centering
\greennote{End of Tutorial – Focus on red flags, traps, and FATF standards.}
\end{document}